%% Tex magic
% !TeX spellcheck = en_US
% !TeX encoding = utf-8
% !TeX root = exercise.tex
%% !TeX program = pdflatex
%% !TeX TXS-program:compile = txs:///pdflatex/[--shell-escape]
%% !BIB program = bibtex

\def\issolution{1}

\NeedsTeXFormat{LaTeX2e}
\documentclass[11pt, a4paper, DIV12, BCOR0mm, headinclude , oneside]{scrartcl}

% Packages
\usepackage[utf8]{inputenc}
\usepackage{booktabs}
\usepackage[english]{babel}
\usepackage{graphicx}
\usepackage{amsmath}
\usepackage{amssymb}
\usepackage{bm}
\usepackage{bbm}
\usepackage{xspace}
\usepackage{url}
\usepackage{fancyhdr}
\usepackage{ifthen}
\usepackage{multicol}
\usepackage{standalone}

\usepackage{enumitem}
\renewcommand{\theenumi}{(\alph{enumi})} % Style of references to enumerate items
\usepackage{color}
\usepackage[colorlinks=false, pdfborder={0 0 0}]{hyperref}
\usepackage{mathtools}
\usepackage{physics}
\usepackage{siunitx}
%\usepackage{minted}
\usepackage{standalone}

% Collections
\usepackage{xcolor}
\definecolor{yellow}{RGB}{255, 198, 0}
\definecolor{orange}{RGB}{255, 130, 0}
\definecolor{blue}{RGB}{0, 32, 91}
\definecolor{red}{RGB}{198, 53, 39}
\definecolor{magenta}{RGB}{138, 27, 97}
\definecolor{lightblue}{RGB}{0, 159, 223}
\definecolor{green}{RGB}{0, 155,119}
\definecolor{lightgreen}{RGB}{132, 189,0}
\definecolor{greenyellow}{RGB}{208, 223,0}

\usepackage{tikz}
\usepackage{tikz-3dplot}

\usetikzlibrary{calc}
\usetikzlibrary{fit}
\usetikzlibrary{backgrounds}
\usetikzlibrary{arrows}
\usetikzlibrary{positioning}
\usetikzlibrary{calc}
\usetikzlibrary{decorations.pathreplacing}
\usetikzlibrary{shapes.multipart}
\usepackage{tikz-3dplot}

% matplotlib2tikz
\usepackage{pgfplots}
\newlength\figureheight
\newlength\figurewidth

% Usage:
% \def\issolution{1}
%
% \begin{solution}
%     $a = 42$
% \end{solution}
%
% If the "issolution" variable is not defined, the solution will not appear.
\usepackage[framemethod=tikz]{mdframed}

\usepackage{comment}
\ifdefined\issolution
    \newenvironment{solution}
    {
    \par\noindent\normalfont\textbf{Solution}\par\nopagebreak%
        \begin{mdframed}[
            linewidth=1pt,
            linecolor=black,
            bottomline=false,
            topline=false,rightline=false,
            innerrightmargin=0pt,
            innertopmargin=0pt,
            innerbottommargin=0pt,
            innerleftmargin=1em,
            skipabove=.5\baselineskip
        ]
    }
    {
        \end{mdframed}
    }
\else
\excludecomment{solution}
\fi


% Commands
\newcommand{\Datum}{\today}
\newcommand{\Vorlesung}{Statistical and Machine Learning}
\newcommand{\CopyrightJahr}{2018}
\newcommand{\EMail}{drude@nt.upb.de}
\newcommand{\AnzahlAnAufgaben}{1}
\newcommand{\mysubtitle}{Exercises}
\newcommand{\thesisauthor}{M.Sc. Lukas Drude}
\newcommand{\thesisdate}{\today}
\newcommand{\fullname}{M.Sc. Lukas Drude}
\newcommand{\university}{Paderborn University}
\newcommand{\department}{Communications Engineering}
\newcommand{\rfootlabel}{Solutions}
\newcommand{\lfootlabel}{Solutions}

\renewcommand{\Pr}{\mathrm{Pr}}
\newcommand{\T}{\mathsf{T}}
\DeclareMathOperator*{\argmax}{argmax}
\DeclareMathOperator*{\argmin}{argmin}

\newcommand{\Uebunghilfe}[1]{{\vfill\hrulefill\\Hint:\\[-2\baselineskip]{#1}}}

\newcommand{\Hinweis}{Achten Sie bei Skizzen auf die korrekte und vollständige Beschriftung der Achsen und markanter Punkte!}

% Header and footer
\lhead{\includegraphics[height=8mm]{latex/upb_logo_rgb_09.pdf}}
\rhead{\tiny\begin{tikzpicture}
\node [fill=blue, text=white, text width=20em, text depth=1em, align=left] (upper) {\textsf{Institut Elektrotechnik und Informationstechnik}};
\node [fill=white, text=black, text depth=-0.1em] (lower) at ($(upper)+(5em,-0.75em)$) {\textbf{\textsf{Nachrichtentechnik}}~~~~\textsf{Prof. Dr.-Ing. Reinhold Häb-Umbach}};
\begin{scope}[on background layer]
\node [fill=blue, fit=(lower), inner sep=0.1em, xshift=0.15em, yshift=-0.15em] (fit) {};
\end{scope}
\end{tikzpicture}% Trailing comment to avoid whitespace.}
\lfoot{Exercise: \Vorlesung, \Datum}
\cfoot{}
\rfoot{\thepage}

% Page layout
\raggedbottom

\renewcommand{\topfraction}{0.9}
\renewcommand{\bottomfraction}{0.9}
\renewcommand{\textfraction}{0.1}

\setcounter{topnumber}{4}
\setcounter{bottomnumber}{0}
\setcounter{totalnumber}{4}

\setlength{\oddsidemargin}{0mm}
\setlength{\evensidemargin}{-5mm}
\setlength{\textwidth}{165mm}
\setlength{\topmargin}{-10mm}
\setlength{\textheight}{235mm}
\setlength{\headheight}{10mm}
\setlength{\parindent}{0pt}

\pagestyle{fancy}


\begin{document}

\newcommand{\rarray}[1]{\renewcommand{\arraystretch}{#1}}
\section{Algorithmic Differentiation}

Algorithmic differentiation allows to differentiate in a complicated computational graph just by combining gradients of individual graph elements.

This exercise aims to teach the difference between forward and backward differentiation and should show, why the computational graph has to be analyzed before backward differentiation is possible.

\tikzset{
    every text node part/.style={align=center},
    >=stealth,
    block/.style={rectangle,draw,minimum height=2em, minimum width=4em,text height=1.5ex,text depth=.25ex},
    branch/.style={circle, fill=black, minimum size=0.25em,inner sep=0},
    arrow/.style={->,shorten >=0.1em},
    reverse arrow/.style={<-,shorten <=0.1em},
    apply/.style={circle, draw, minimum size=0.8em, inner sep=0},
    plus/.style={apply, label=center:{+}},
    times/.style={apply, label=center:{$\times$}}
}

\begin{center}
\begin{tikzpicture}[node distance=3em and 4em]
\node (mse) [block] {MSE};
\draw [arrow] (mse) -- +(0, 3em) node [below right] {$J$};
\node (y_hat) [plus, below=of mse] {};
\draw [arrow] (y_hat) -- node [right] {$\hat y$} (mse);
\node (relu) [block, below=of y_hat] {ReLU};
\draw [arrow] (relu) -- node [right] {$h$} (y_hat);
\draw [reverse arrow] (mse) -- +(4em, 0) node [above left] {$y$};
\node (z) [branch, below=of relu] {};
\draw [arrow] (z) -- node [right] {$z$} (relu);
\draw [arrow] (z) -- +(4em, 0) |- (y_hat);
\node (plus) [plus, below=of z] {};
\draw (plus) -- (z);
\draw [reverse arrow] (plus) -- +(4em, 0) node [above left] {$b$};
\node (times) [times, below=of plus] {};
\draw [arrow] (times) -- node [right] {$v$} (plus);
\draw [reverse arrow] (times) -- +(4em, 0) node [above left] {$w$};
\draw [reverse arrow] (times) -- +(0, -3em) node [above right] {$x$};
\end{tikzpicture}
\end{center}

\begin{enumerate}
\item Obtain a valid execution order of the operations in the displayed graph.
Start with the cost $J$ and determine, which values are necessary to apply the current operation, e.g. $\mathrm{MSE}(\cdot, \cdot)$.
\item Calculate each value in the graph.
Use the seed values $y=7$, $b=-3$, $w=2$, and $x=5$.
\item Derive the symbolic partial derivatives of each operation independently.
\item Calculate the numeric forward gradients for each graph element using the following seeds:
\begin{tabular}{
l
@{\hskip 0ex)\hskip 1ex}
c
@{\,$=1$ others 0}
}
a & $\dot x_{\mathrm{seed}}$ \\
b & $\dot w_{\mathrm{seed}}$ \\
c & $\dot b_{\mathrm{seed}}$ \\
d & $\dot y_{\mathrm{seed}}$ \\
\end{tabular}
\item Calculate the numeric backward gradient with the seed $\bar J_{\mathrm{seed}}=1$.
\end{enumerate}

\Uebunghilfe{
\begin{align*}
\mathrm{heaviside}(z) &= \begin{cases}
0, & z < 0 \\
1, & z \geq 0
\end{cases}
\end{align*}
}

\clearpage
\begin{solution}
\begin{center}
\rarray{1.3}
\begin{tabular}{lcc}
\toprule
Operation & Value & Symbolic gradients \\
\midrule
$x$ & 5 \\
$w$ & 2 \\
$b$ & -3 \\
$y$ & 7 \\
$v=wx$ & 10 & $\pdv{v}{w} = x$, $\pdv{v}{x} = w$ \\
$z=v+b$ & 7 & $\pdv{z}{v} = 1$, $\pdv{z}{b} = 1$\\
$h=\mathrm{ReLU}(z)$ & 7 & $\pdv{h}{z}=\mathrm{heaviside}(z)$ \\
$\hat y = h + z$ & 14 & $\pdv{y}{h}=1$, $\pdv{y}{z}=1$ \\
$J=\mathrm{MSE}(\hat y, y)$ & 49 & $\pdv{J}{\hat y}=2(\hat y - y)$, $\pdv{J}{y}=-2(\hat y - y)$ \\
\bottomrule
\end{tabular}
\end{center}

\begin{center}
\rarray{1.3}
\begin{tabular}{lccccc}
\toprule
Operation & \multicolumn{4}{c}{Forward mode} & Backward mode \\
& $\dot b_{\mathrm{seed}}=1$ & $\dot w_{\mathrm{seed}}=1$ & $\dot x_{\mathrm{seed}}=1$ & $\dot y_{\mathrm{seed}}=1$ & $\bar J_{\mathrm{seed}}=1$ \\
\midrule
$x$ & $\dot x = 0$ & $\dot x = 0$ & $\dot x = 1$ & $\dot x = 0$ & $\color{blue}\bar x = 56$ \\
$w$ & $\dot w = 0$ & $\dot w = 1$ & $\dot w = 0$ & $\dot w = 0$ & $\color{green}\bar w = 140$ \\
$b$ & $\dot b = 1$ & $\dot b = 0$ & $\dot b = 0$ & $\dot b = 0$ & $\color{red}\bar b = 28$ \\
$y$ & $\dot y = 0$ & $\dot y = 0$ & $\dot y = 0$ & $\dot y = 1$ & $\color{cyan}\bar y = -14$ \\
$v=wx$ & $\dot v = 0$ & $\dot v = 5$ & $\dot v = 2$ & $\dot v = 0$ & $\bar v = 28$ \\
$z=v+b$ & $\dot z = 1$ & $\dot z = 5$ & $\dot z = 2$ & $\dot z = 0$ & $\bar z = 28$ \\
$h=\mathrm{ReLU}(z)$ & $\dot h = 1$ & $\dot h = 5$ & $\dot h = 2$ & $\dot h = 0$ & $\bar h = 14$ \\
$\hat y = h + z$ & $\dot y = 2$ & $\dot y = 10$ & $\dot y = 4$ & $\dot y = 0$  & $\bar v = 14$\\
$J=\mathrm{MSE}(\hat y, y)$ & $\color{red}\dot J = 28$ & $\color{green}\dot J = 140$ & $\color{blue}\dot J = 56$ & $\color{cyan}\dot J = -14$ & $\bar J = 1$ \\
\bottomrule
\end{tabular}
\end{center}

Example:
\begin{align*}
\dot v &= \pdv v b = \pdv v w \pdv wb + \pdv v x \pdv x b \\
&= \pdv v w \dot w + \pdv v x \dot x
\end{align*}


\begin{align*}
\bar v &= \pdv J v = \dots = \bar z \pdv z v
\end{align*}


\end{solution}

\end{document}
